\section*{Problem and Motivation}

Collision detection in real time with a large number of moving objects is challenging. Comparing
each object with each of the others is a possible solution. However, this brute-force approach requires a
time complexity of $O({n}^{2})$, where $n$ is the number of objects. The simulation runs become slower with an
increase in the number of objects. Applications like physics and video games that need many spatial
queries can be severely limited.

This project uses spatial partitioning to speed up collision detection. In particular, this project
will use and measure the performance of Quadtree. A Quadtree is a data structure used for recursive
subdivision of two-dimensional space. The project will measure and compare the performance of
Quadtree against a brute-force approach and a uniform spatial grid.

An interesting problem of using spatial partitioning is how different spatial partitioning
structures perform depending on the particle distribution. Although Quadtree can reduce unnecessary
checks, a lot of overhead is introduced by having a structure. The simulation particle count and
distribution as well as Quadtree parameters need to be examined.
